\documentclass[10pt,a4paper]{amsart}
\usepackage[margin=19mm]{geometry}
\usepackage{amsmath,amssymb,mathtools}
\usepackage[scr=rsfs,cal=euler]{mathalfa}
\usepackage{xfrac}
\usepackage[unicode,hidelinks,bookmarks=false]{hyperref}
\newcommand{\ie}{i.e.\ }
\newcommand{\cf}{cf.\ }
\newcommand{\resp}{respectively\ }
\newcommand{\Subsection}{\subsection}
\providecommand{\nobreakdash}{\nobreak\hbox{-}}
\setlength{\emergencystretch}{2em}
\multlinegap0pt
\usepackage{amssymb}
\usepackage{mathtools}
\usepackage{float}
\usepackage{enumerate}
\newtheorem{corollary}{Corollary}
\newtheorem{proposition}{Proposition}
\newcommand{\E}{\mathrm{E}}
\renewcommand{\H}{\mathrm{H}}
\renewcommand{\L}{\mathrm{L}}
\newcommand{\R}{\mathbb{R}}
\newcommand{\T}{\mathbb{T}}
\newcommand{\cT}{\mathcal{T}}
\newcommand{\p}{\mathrm{p}}
\title{The topology of tridiagonal isospectral sets}
\author{David Mart\'inez Torres}
\date{Source: arXiv:2609.05733v1 (CC BY 4.0)}
\begin{document}
\maketitle

\noindent\emph{Source and licence.} The topology of tridiagonal isospectral sets. Version arXiv:2609.05733v1 (CC BY 4.0), distributed by arXiv under the Creative Commons Attribution 4.0 International licence, http://creativecommons.org/licenses/by/4.0/, as recorded in the arXiv licence metadata for this exact version and shown on the abstract page https://arxiv.org/abs/2609.05733v1. Full licence notice: \texttt{licenses/arxiv-2609.05733-CC-BY-4.0.txt}.\par\smallskip

\noindent\textit{Editorial scope.} Counted unit: the article's proposition pro:contractible (source0.tex lines 673--674). The article's complete original proof of it is delivered with it (source0.tex lines 675--703, one \texttt{proof} environment of 29 lines). No mathematics is written, repaired or re-derived: the statement and the proof are the article's own lines, in the article's own notation, and no retained line is substituted or re-flowed.  Retained uncounted as the source-local context the proof needs: lines 571--581.  Nothing was omitted from any retained span, so the package carries no omission record.  No retained line contains a citation, so the wrapper carries no bibliography and no bibliography entry.  No retained line contains a figure environment, an image-inclusion command or drawing code, so no image is delivered and the package declares no asset dependency.  Editorial formatting only, in addition: an AMS \texttt{amsart} preamble that restates the article's own declarations and macros used by the retained lines, namely the environments \texttt{corollary}, \texttt{proposition} on separate counters, together with the standard packages those lines need.  The retained environments print in this document as Corollary 1, Proposition 1 in order of appearance, whereas the article prints the same items with the numbering of its own full document.  The complete native source of the article as distributed in the arXiv e-print for this exact version is bundled unchanged as \texttt{sources/209462/source0.tex}.
\begin{corollary}\label{cor:morse-ordering} The function $f:\cT\to \R$
\begin{enumerate}
 \item does not have local maxima and minima different from the global ones;
 \item if $\Lambda$ is the maximum (resp. minimum) for $f$ on $\cT_{\p}$, then it cannot be the maximum (resp.
 minimum) for $f$ on any other submanifold component through $\Lambda$. Moreover,
 \begin{itemize}
  \item $\p_M\in \mathfrak{M}(\Lambda)$ (resp. $\p_m$) is unique and equals $\p$;
  \item $W^s(\p_M)$ (resp. $W^u(\p_m)$) strictly contains the  stable (resp. unstable) manifolds for all the other $\cT_{\p'}$, $\p'\in \mathfrak{M}(\Lambda)$.
 \end{itemize}
\end{enumerate}
\end{corollary}

\begin{proposition}\label{pro:contractible} The positive chamber $\Delta$ is contractible.
 \end{proposition}

\begin{proof}
We contract $\Delta$ to the minimum $\Lambda^-$ by adapting the
standard homotopy $h_\Lambda$ to the positive chamber.

%
% We go back to
% the auxiliary coordinate space $\L^\Lambda_\H$ and the diagonal linear vector field $Y$.
% The group $\E$ acts on $\L^\Lambda_\H$ and the composition of the inclusion $\L^\Lambda_\H$ with $\Psi$ is $\E$ equivariant. The action on $\L^\Lambda_\H$ amounts to sign changes the subdiagonal variables.
We denote by $\Delta^\Lambda\subset \T^\Lambda$ the union of positive chambers for the action of $\E$ on $\T_\p$, $\p\in \mathfrak{M}(\Lambda)$. The equivariance of $\Psi:\T^\Lambda\to \cT$ implies that $\Delta^\Lambda$ is mapped (in)to $\Delta$. The positive chambers $\Delta$ and $\Delta^\Lambda$ are invariant by the Toda vector field and by $Y$, respectively; $\Delta^\Lambda$ is also invariant by $\sum_l -U_l\frac{\partial}{\partial U_l}$. Therefore we can restrict the standard homotopy

\begin{equation}\label{eq:fundamental-standard-homotopy}
  h_\Lambda:f^{-1}((-\infty, f(\Lambda)+2\epsilon])\cap \Delta^\Lambda\to \left(f^{-1}((-\infty, f(\Lambda)-\epsilon])\cup e_{\p_M}\right)\cap \Delta^\Lambda.
\end{equation}
The intersection $W^s(\p_M)\cap \Delta^\Lambda$ is a positive chamber for the residual action of $\mathrm{E}$ on $W^s(\p_M)$. It has positive dimension  unless $W^s(\p_M)$ is a point. By Corollary \ref{cor:morse-ordering} this only happens at $\Lambda^-$. If this is not the case, then we consider the auxiliary constant diagonal vector field
\[\sum_{k}\epsilon_k\frac{\partial}{\partial S_k},\]
where $\epsilon_k=0$ if $S_k$ is not a coordinate in $W^s(\p_M)$, and it is 1 or $-1$ otherwise, so that it points inside the positive chamber $W^s(\p_M)\cap \Delta^\Lambda$ along the part of its boundary which is a union of coordinate subspaces. We define the homotopy
\[h_3:\left(f^{-1}((-\infty, f(\Lambda)-\epsilon])\cup e_{\p_M}\right)\cap \Delta^\Lambda\to f^{-1}((-\infty, f(\Lambda)-\epsilon])\cap \Delta^\Lambda\]
that is the identity  at a point in $f^{-1}((-\infty, f(\Lambda)-\epsilon])\cap \Delta^\Lambda$, and that at  $L\in e_{\p_M}\cap \Delta^\Lambda$ traverses at a constant speed the segment tangent to the constant auxiliary vector field that goes from $L$ to the component of the boundary of  $e_{\p_M}\cap \Delta^\Lambda$ in $\partial e_{\p_M}$.

We denote the concatenation of $h_\Lambda$ with $h_3$ by $h_{\Delta^\Lambda}$. Observe that $h_{\Delta^\Lambda}$ is well-defined in $\Delta$, since $h_3$ takes $\cT_{\p_M}\cap \Delta$ to itself.



We use  the Toda vector field combined with the homotopy $h_{\Delta^\Lambda}$
near each critical point to contract $\Delta$ into $\Lambda^{-}$.
We start at the maximum $\Lambda^+$.
There, we apply $h_3$ to retract $\Delta$ onto $f^{-1}((f(\Lambda^+)-\epsilon))\cap \Delta$. Then we follow the flow of the Toda vector field until we hit the level set
$f^{-1}(f(\Lambda)+2\epsilon))\cap \Delta$, where $\Lambda$ is the next critical point. If $\Lambda$ is not the minimum, then  we can apply  $h_{\Delta^\Lambda}$ to retract onto $f^{-1}((f(\Lambda)-\epsilon))\cap \Delta$. By induction, we retract $\cT$ onto $\Lambda^{-}$.
\end{proof}

\end{document}
